% Lapisan solusi O001 -- Bab 13 (Konstruksi GNS)
% 1 solusi asli terpisah; bukan tulisan John M. Erdman.
% Provenans model: OpenAI Codex gpt-5.6-sol, Ultra, atas arahan pengguna.
% Lisensi komponen solusi: CC BY-SA 4.0.
% Tidak menyiratkan dukungan John M. Erdman atau Portland State University.

\markboth{SOLUSI O001 UNTUK BAB 13: KONSTRUKSI GNS}{SOLUSI O001 UNTUK BAB 13: KONSTRUKSI GNS}
\section*{Solusi O001 untuk Bab 13: Konstruksi GNS}
\addcontentsline{toc}{section}{Solusi O001 untuk Bab 13: Konstruksi GNS}
\begin{center}
\small
Komponen solusi asli terpisah, CC BY-SA 4.0. Ditulis dengan bantuan
OpenAI Codex gpt-5.6-sol, Ultra, atas arahan pengguna. Pernyataan latihan
berasal dari edisi terjemahan karya John M. Erdman; solusi ini bukan tulisan
beliau dan tidak menyiratkan dukungan beliau atau Portland State University.
\end{center}

% O001-SOLUTION-ID: O001-FAOA-2015-CH13-EX-001-SOLUTION
% SOURCE-EXERCISE-ID: FAOA-2015-CH13-NODE-0018
% STATEMENT-TARGET-SHA256: 84c22fa77eaccb119f0537a7807cc91265e0dc1e0c2aae2137eadc8497924228
\begin{o001solution}
  {O001-FAOA-2015-CH13-EX-001-SOLUTION}
  {FAOA-2015-CH13-NODE-0018}
  {84c22fa77eaccb119f0537a7807cc91265e0dc1e0c2aae2137eadc8497924228}
\begin{o001statement}
Misalkan $X$ suatu ruang Hausdorff kompak lokal. Carilah suatu representasi
isometrik (dan karena itu setia) $(\pi,H)$ dari aljabar-$C^*$ $C_0(X)$ pada suatu
ruang Hilbert~$H$.
\end{o001statement}
\begin{o001answer}
Ambil $H=\ell^2(X)$ dan definisikan
$(\pi(f)\xi)(x)=f(x)\xi(x)$. Maka $\pi$ adalah representasi isometrik dan
setia dari $C_0(X)$ pada~$H$.
\end{o001answer}
\begin{o001proof}
Untuk setiap $f\in C_0(X)$ dan $\xi\in\ell^2(X)$, perkalian titik demi titik
memberi
\[
 \|\pi(f)\xi\|_2^2
 =\sum_{x\in X}|f(x)|^2|\xi(x)|^2
 \leq\|f\|_\infty^2\|\xi\|_2^2.
\]
Jadi $\pi(f)$ operator terbatas dan $\|\pi(f)\|\leq\|f\|_\infty$. Definisi
titik demi titik langsung memberi
\[
 \pi(fg)=\pi(f)\pi(g),
 \qquad
 \pi(\overline f)=\pi(f)^*,
\]
serta linearitas~$\pi$.

Untuk setiap $x\in X$, misalkan $e_x$ adalah vektor basis satuan yang bernilai
$1$ di~$x$ dan nol di tempat lain. Maka
\[
 \|\pi(f)e_x\|_2=|f(x)|.
\]
Akibatnya $\|\pi(f)\|\geq |f(x)|$ untuk setiap~$x$, dan dengan mengambil
supremum diperoleh $\|\pi(f)\|\geq\|f\|_\infty$. Bersama estimasi sebelumnya,
$\|\pi(f)\|=\|f\|_\infty$. Jadi $\pi$ isometrik; khususnya
$\pi(f)=0$ memaksa $f=0$, sehingga representasi itu setia. Konstruksi ini
berlaku tanpa perlu memilih ukuran pada~$X$, bahkan ketika $X$ tidak dapat
dihitung.
\end{o001proof}
\end{o001solution}
